\section{Aprendizaje por refuerzo multitarea}

\subsection{Motivación: por qué el RL tarde o temprano tiene que avanzar hacia una generalist policy}

El expositor plantea la motivación del RL multitarea de forma muy directa: los sistemas inteligentes del mundo real casi nunca enfrentan una sola tarea. Un asistente general tiene que encargarse de reservar viajes, hacer compras y responder preguntas; un robot tiene que caminar, agacharse, transportar objetos y limpiar; un sistema de recomendación tiene que atender las preferencias de distintos usuarios. La verdadera pregunta no es «si cierta política puede aprender una tarea», sino «si podemos entrenar una política capaz de compartir conocimiento entre muchas tareas».

\begin{figure}[H]
\centering
\includegraphics[width=0.9\textwidth]{slides-images/slide-11.jpg}
\caption{Panorama de aplicaciones del RL multitarea: desde asistentes basados en LLM hasta robots de manipulación móvil y sistemas de recomendación}
\end{figure}

\begin{figure}[H]
\centering
\includegraphics[width=0.9\textwidth]{slides-images/slide-12.jpg}
\caption{La motivación de fondo del RL multitarea: repartir la complejidad de datos entre las tareas para construir sistemas más general}
\end{figure}

\subsection{Una tarea es, en realidad, un MDP distinto}

\begin{importantbox}{Formalización del MDP multitarea}
El problema multitarea se escribe como un conjunto de MDP $\{\mathcal{M}_i\}$; cada tarea puede variar en el espacio de estados, el espacio de acciones, la distribución del estado inicial, la dinámica y la función de recompensa. En el entrenamiento, es habitual incorporar el identificador de la tarea $z_i$ al estado, lo que se escribe como:
$$\tilde{\mathcal{M}} = (\mathcal{S} \times \mathcal{Z}, \mathcal{A}, \tilde{p}, \tilde{r}, \gamma)$$
Así, la política se convierte en una política condicionada $\pi(a\mid s, z)$.
\end{importantbox}

Esta definición es más amplia que muchas nociones intuitivas de «etiqueta de tarea». Una tarea no consiste necesariamente solo en «hacer cosas distintas»; también puede tratarse de usuarios distintos, configuraciones iniciales distintas, ajustes de dinámica distintos o incluso estructuras corporales distintas.

\begin{figure}[H]
\centering
\includegraphics[width=0.9\textwidth]{slides-images/slide-14.jpg}
\caption{Ejemplos de distribuciones de tareas: las diferencias entre tareas en sistemas de recomendación, animación de personajes y RL con múltiples robots}
\end{figure}

\subsection{Cómo se le proporciona a la política el identificador de la tarea}

La descripción de la tarea $z$ puede adoptar varias formas:
\begin{itemize}
    \item \textbf{One-hot task ID}: es sencillo de implementar, pero casi no permite la generalización sin ejemplos previos.
    \item \textbf{Instrucciones en lenguaje natural}: son más adecuadas para tareas de vocabulario abierto y se integran con mayor facilidad con VLA/LLM.
    \item \textbf{Estado objetivo}: expresa la tarea como «a dónde ir».
    \item \textbf{Parámetros de recompensa o demostraciones en video}: convienen cuando la tarea tiene una semántica más compleja.
\end{itemize}

\begin{figure}[H]
\centering
\includegraphics[width=0.9\textwidth]{slides-images/slide-15.jpg}
\caption{El identificador de la tarea puede ser un índice, lenguaje, una demostración en video o un estado objetivo}
\end{figure}

\begin{figure}[H]
\centering
\includegraphics[width=0.9\textwidth]{slides-images/slide-16.jpg}
\caption{Al incorporar el identificador de la tarea al estado, el RL multitarea puede reescribirse como un MDP estándar}
\end{figure}

\begin{knowledgebox}{La importancia de esta reescritura}
Una vez que se acepta la notación $s=(\bar{s}, z)$, el RL multitarea ya no necesita un aparato matemático de RL completamente nuevo. La policy, la Q-function y el critic del caso de una sola tarea pueden reescribirse casi directamente en forma condicionada:
\begin{itemize}
    \item $\pi_\theta(a\mid s)\rightarrow \pi_\theta(a\mid \bar{s}, z)$
    \item $Q_\phi(s,a)\rightarrow Q_\phi(\bar{s}, a, z)$
\end{itemize}
Por esta razón, la implementación del RL multitarea resulta mucho menos ajena de lo que parece a primera vista.
\end{knowledgebox}

\subsection{Cómo se lleva a cabo el imitation learning multitarea}

La forma más directa de aprendizaje multitarea sigue siendo la clonación de comportamiento: se reúnen datos de demostración de varias tareas y se entrena una política condicionada $\pi_\theta(a\mid s,z)$. Sin embargo, el curso señala de manera especial una técnica práctica que suele pasarse por alto: \textbf{stratified sampling}. Es decir, procurar que cada mini-batch contenga datos de varias tareas, para evitar que las tareas más frecuentes dominen el entrenamiento.

\begin{figure}[H]
\centering
\includegraphics[width=0.9\textwidth]{slides-images/slide-18.jpg}
\caption{Imitation learning multitarea: una política condicionada combinada con muestreo estratificado reduce la varianza del gradiente}
\end{figure}

\begin{knowledgebox}{Por qué el aprendizaje multitarea puede compartir datos}
Las distintas tareas suelen compartir subestructuras de habilidades. Por ejemplo, «tomar el cubo rojo» y «tomar el cubo azul» comparten la acción de sujetar; «limpiar la mesa» y «ordenar la cocina» comparten la capacidad de desplazarse y ubicarse. El valor esencial del aprendizaje multitarea consiste en que estas habilidades locales se reutilicen dentro de parámetros compartidos, en lugar de aprenderse desde cero por separado para cada tarea.
\end{knowledgebox}

\subsection{De BC-Z a OpenVLA: formas actuales de introducir el condicionamiento por tarea}

Las slides también presentan expresamente dos tipos de ejemplos que muestran cómo el task conditioning se ha fusionado en los últimos años con los foundation model:
\begin{itemize}
    \item Los métodos del tipo de BC-Z usan un task embedding explícito o condiciones multimodales, con el objetivo de lograr la generalización zero-shot entre tareas.
    \item Sistemas como OpenVLA son más directos: introducen la descripción de la tarea como prompt en un modelo de visión-lenguaje-acción.
\end{itemize}

\begin{figure}[H]
\centering
\includegraphics[width=0.9\textwidth]{slides-images/slide-19.jpg}
\caption{BC-Z: generalización sin ejemplos previos entre tareas mediante el condicionamiento por tarea}
\end{figure}

\begin{figure}[H]
\centering
\includegraphics[width=0.9\textwidth]{slides-images/slide-20.jpg}
\caption{OpenVLA: la descripción de la tarea se introduce como prompt en el modelo VLA}
\end{figure}

\begin{figure}[H]
\centering
\includegraphics[width=0.9\textwidth]{slides-images/slide-21.jpg}
\caption{Ejemplo robótico de imitation multitarea: de «doblar la toalla» a «colgar la toalla»}
\end{figure}

\begin{figure}[H]
\centering
\includegraphics[width=0.9\textwidth]{slides-images/slide-22.jpg}
\caption{De unas cuantas tareas a tareas domésticas abiertas: la escala de las políticas multitarea crece con rapidez}
\end{figure}

\subsection{Resumen de la sección}

Lo esencial del RL multitarea no es «meter más tareas en una sola red», sino expresar, mediante el condicionamiento, las distintas tareas como un conjunto de problemas relacionados que una misma política debe resolver. Siempre que exista una estructura compartida entre las tareas, el entrenamiento multitarea puede repartir la complejidad de datos.
